\section{Introduction}

A practitioner applying CCNet's standard perplexity quality filter to the RedPajama-V2 dataset would retain 86\% of Spanish documents but only 16\% of English documents --- not because English web text is lower quality, but because a single global percentile threshold applied to cross-language perplexity scores conflates quality with language identity.

This outcome is not a corner case. At $k$=40 (retaining the 40th percentile of documents with lowest perplexity), virtually all Spanish documents pass while 79\% of German documents are excluded. Cram{\'e}r's $V$ --- a standard measure of association magnitude in contingency tables --- reaches 0.57 at $k$=40, placing the language $\times$ retention relationship firmly in the ``large'' effect range by Cohen's conventions. The disparity is structural: the ordering of retention rates (\textit{es} $>$ \textit{fr} $>$ \textit{it} $>$ \textit{en} $>$ \textit{de}) maps exactly onto the Romance--Germanic language family divide and holds across every threshold level we test.

The surface problem is familiar: multilingual pretraining datasets require quality filtering to remove low-quality web text, and perplexity-based filtering using language models trained on curated reference corpora has become a standard approach, popularized by CCNet \cite{Wenzek2019CCNet} and widely adopted in datasets including RedPajama-V2 \cite{Weber2024RedPajama}. The deeper problem is less recognized: when CCNet-style perplexity scores are computed per-language but combined under a single global threshold, the cross-language incomparability of those scores creates a retention disparity that tracks language family membership rather than document quality. CCNet's original per-language design calibrates thresholds separately per language via language-specific cutoff values; practitioners adapting this pipeline for large-scale datasets often replace per-language cutoffs with a single global $k$-th percentile, losing the calibration that CCNet intended. The gap we address is quantitative: no prior work provides a calibrated baseline measurement of how large this disparity is, using a standard associativity metric across multiple threshold levels on a real multilingual dataset.

Our key insight is that CCNet's per-language KenLM language models, trained on language-specific Wikipedia corpora of structurally different sizes and domain compositions, produce perplexity scales that are incomparable across language families. A global $k$-th percentile threshold is then anchored by the combined distribution, which is dominated by Germanic-language documents clustering at lower absolute perplexity values. The result is not a quality filter but a language-family membership test.

This paper makes four contributions, each building directly on this insight:

\begin{enumerate}
\item \textbf{Calibrated baseline measurement.} We provide the first precise, reproducible measurement of language-group retention disparity (Cram{\'e}r's $V = 0.40$--$0.57$) from global \texttt{ccnet\_perplexity} thresholding on RedPajama-V2 ($n = 208,262$ documents, 5 languages), with Holm-corrected $p$-values $\approx 0$ for all five threshold levels $k \in \{10, 20, 30, 40, 50\}$.

\item \textbf{Cross-language family structure characterization.} We show that the disparity is not random --- the retention ordering \textit{es} $>$ \textit{fr} $>$ \textit{it} $>$ \textit{en} $>$ \textit{de} is consistent across all five threshold levels, providing structural evidence consistent with the CCNet KenLM mechanism.

\item \textbf{Prior estimate recalibration.} Literature-derived estimates predict Cram{\'e}r's $V \in [0.29, 0.41]$; our empirical measurement yields $V = 0.40$--$0.57$ --- 25--40\% larger. Researchers sizing correction studies from CCNet descriptions will systematically underestimate the required effect size.

\item \textbf{Reusable statistical methodology.} We demonstrate a CPU-only, static-file pipeline using pre-computed \texttt{ccnet\_perplexity} quality signals from RedPajama-V2, Cram{\'e}r's $V$ as the association measure, and Holm-Bonferroni correction --- a methodology replicable without GPU infrastructure or model retraining.
\end{enumerate}

The remainder of the paper is organized as follows. Section~\ref{sec:related} positions our work relative to multilingual quality filtering and web corpus curation literature. Section~\ref{sec:methodology} describes the measurement methodology. Section~\ref{sec:setup} presents the experimental setup. Section~\ref{sec:results} reports results. Section~\ref{sec:discussion} discusses implications and limitations. Section~\ref{sec:conclusion} concludes with a call for per-language calibration as a default practice in multilingual dataset curation.
